\documentclass{article}

\usepackage{amsmath, parskip, tikz, hyperref}
\usetikzlibrary{automata}
\usetikzlibrary{positioning}
\usetikzlibrary{arrows}
\usepackage[margin=1.25in]{geometry}
\usepackage[shortlabels]{enumitem}

\newcommand\blanksymbol\textvisiblespace

\renewcommand{\thesection}{\arabic{section}.}
\renewcommand{\thesubsection}{\alph{section}.}

\hypersetup{
    colorlinks=true,
    urlcolor=blue,
    linkcolor=blue
}

\tikzset{
    node distance=2.5cm,
    every state/.append style={
        semithick,
        fill=gray!10
    },
    initial/.append style={
        initial text={},
        initial distance=0.5cm
    },
    accepting/.append style={
        double=gray!10,
        double distance=2pt,
        outer sep=1pt
    },
    every edge/.append style={
        draw,
        ->,>=stealth',
        auto,
        semithick
    }
}

\title{Homework 7: Undecidability}
\author{Will Griffin}

\begin{document}
    \maketitle

    \section{Bounds checking}

    \begin{enumerate}[(a)]
        \item Prove that it is undecidable whether a TM $M$, on input $w$, ever attempts to move its head past the left end of the tape.

            Let $\text{LEFT}_\text{TM}$ be the langauge where 

            \begin{equation*}
            \text{LEFT}_\text{TM} = \{ \langle M, w \rangle \mid M \text { attempts to move its head past the left end of the tape}\}
            \end{equation*}

            Assume $\text{LEFT}_\text{TM}$ is decidable by a TM $R$.

            Construct a TM $S$ that decides $\text{A}_\text{TM}$.

            $S = $ ``On input $\langle M, w \rangle$,
            \begin{enumerate}[1.]
                \item Construct a TM $M' = $ ``On input x:
                    \begin{enumerate}[a.]
                        \item Simulate $M$ on $\texttt{\#}w$. If the simulation of $M$ ever tries to move to the cell containing \texttt{\#}, move the head to the right and continue the simulation.

                        \item If $M$ halts, attempt to move the head of $M'$ left past the first cell."
                    \end{enumerate}

                \item Run $R$ on $\langle M', \varepsilon \rangle$.

                \item If $R$ accepts, accept.

                \item If $R$ rejects, reject."
            \end{enumerate}

            On a case basis:
            \begin{enumerate}[-]
                \item $M$ accepts $w$: $M'$ will move left past the first cell. $R$ accepts $\langle M', \varepsilon \rangle$. $S$ accepts $\langle M, w \rangle$.

                \item $M$ rejects $w$: $M'$ will move left past the first cell. $R$ accepts $\langle M', \varepsilon \rangle$. $S$ accepts $\langle M, w \rangle$.

                \item $M$ loops on $w$: $M'$ will not move left past the first cell. $R$ rejects $\langle M', \varepsilon \rangle$. $S$ rejects $\langle M, w \rangle$.
            \end{enumerate}

            Since $S$ decides $\text{A}_\text{TM}$, but $\text{A}_\text{TM}$ is undecidable, we have a contradiction. Therefore, $\text{LEFT}_\text{TM}$ is undecidable.

        \item Prove that it is decidable whether a TM $M$, on input $w$, ever attempts to move its head past the right end of the input string $w$. Your answer should be a high-level description of a TM. 

            Construct a TM $R$ that decides whether a TM $M$, on input $w$, attempts to move its head past the right end of $w$. This is done by simulating $M$ on $w$. The total number of configurations $M$ can be in, $C$, is equal to $C = |Q| \times |w| \times |\Gamma|^{|w|}$. Simulate $M$ on $w$ for at most $C$ steps. Between each step in the situation, if the head moves right past the position $n$, \textit{accept}. If $M$ halts without ever moving right past the position $n$, \textit{reject}. After $C$ steps of simulation, if $M$ has not halted, \textit{reject}.
            
    \end{enumerate}

    \section{The Power of 10}

    Know:
    \begin{enumerate}[-]
        \item A 10-compliant TM always halts on every input
            
        \item It is decidable whether a TM is 10-compliant
    \end{enumerate}

    Consider the language

    \begin{equation*}
        L_2 = \{\langle M \rangle \mid M \text{ is a 10-compliant TM that rejects } \langle M \rangle\}
    \end{equation*}

    \begin{enumerate}[(a)]
        \item Prove that $L_2$ is decidable.

            If it is decidable whether a TM $M$ is 10-compliant, then there is some decider $R$ that decides 10-compliancy for a TM $M$.

            To prove that $L_2$ is decidable, construct a TM $S$ that operates on some input $w$. If $w \ne \langle M \rangle$ ($w$ is not a valid TM encoding), \textit{reject}. Run the decider $R$ on $\langle M \rangle$. If $M$ is not 10-compliant, \textit{reject}. If $M$ is 10-compliant, simulate $M$ on $\langle M \rangle$. Since $M$ is 10-compliant, this simulation will eventually halt. If $M$ rejects $\langle M \rangle$, \textit{accept}, otherwise, if $M$ accepts $\langle M \rangle$, \textit{reject}.

            Since this TM $S$ halts on any input and only accepts TMs that are both 10-compliant and reject their own encoding, $L_2$ is decidable.

        \item Prove that any TM that decides $L_2$ must not be 10-compliant.

            Any TM whose encoding is in $L_2$ must reject its own encoding. However, $S$, the decider for $L_2$, accepts if the TM it operates on rejects itself. Therefore, for $S$ to be 10-compliant, it must reject itself, but by determining whether $S$ is in $L_2$, if $S$ rejects $\langle S \rangle$, then $S$ checking $\langle S \rangle$ accepts. This causes a paradox.

            By this contradiction, any decider for $L_2$ is not 10-compliant.

        \item Where in your solution to (a) is the violation of 10-compliant? (Full credit for any reasonable answer; I just want you to think about it).

            Since deciding $L_2$ is a non-10-compliant operation, any machine that decides $L_2$ is not 10-compliant. Therefore, in running $S$ on $\langle S \rangle$, the 10-compliancy-decider step in $S$ breaks 10-compliancy for $S$.

    \end{enumerate}

    \section{Grep and Sed} 

    \begin{enumerate}[(a)]
        \item Prove that it is decidable whether 2 regular expressions are equivalent to each other. Hint: prove that it is decidable whether a DFA recognizes the empty language. Then, use closure properties.

            For two regular expressions $\alpha_1$ and $\alpha_2$ that recognize the languages $L_1$ and $L_2$ respectively, if $\alpha_1$ and $\alpha_2$ are equivalent, then $L_1 \cap L_2 = L_1 \cup L_2$. The difference between these two languages ($L_\text{diff} = (L_1 \cap \overline{L_2}) \cup (L_2 \cap \overline{L_1})$) would then be the empty language $L_\text{diff} = \emptyset$.

            Therefore, since the closure operations of complementation, union, and intersection are closed under regular languages, it is enough to prove whether determining a DFA recognizes the empty language $L_\text{diff}$ is decidable.

            To do this, construct some TM $M_\emptyset$ that operates on the input string of $\langle M_\text{diff} \rangle$ (the DFA constructed from $L_\text{diff}$). Search for a path to any accept state in $M_\text{diff}$. If any accept state is reachable, \textit{reject}. If no accept state is reachable, \textit{accept}.

            Since a TM that halts on any input of $\langle M_\text{diff} \rangle$ can be constructed to decide $L_\text{diff}$, $L_\text{diff}$ is decidable. By closure properties of complementation, union, and intersection under regular lagnuages, this can be extended to say that it is decidable whether 2 regular expressions are equivalent.

        \item Prove that it is undecidable whether two \texttt{sed} programs (as defined in CP3) are equivalent to each other. Hint: use CP3 part 3

            In CP3 part 3, we created a program that can convert any Turing machine into an equivalent \texttt{sed} script. By this implementation, this proves \texttt{sed} to be \href{https://en.wikipedia.org/wiki/Turing_completeness}{Turing-complete}. Turing-completeness means that the model of computation is at least equal in power to Turing machines. However, since Turing machines are the most powerful method of computation, \texttt{sed} is equal in computational power to Turing machines. Therefore, we can say that for two \texttt{sed} programs $S_1$ and $S_2$, we can construct two equivalent TMs $M_1$ and $M_2$ respectively. However, by Theorems 5.4 and 5.2 (Sipser), is is undecidable whether two TMs are equivalent. Therefore, because \texttt{sed} is Turing-complete, it is also undecidable whether two \texttt{sed} programs are equivalent to each other.

    \end{enumerate}
\end{document}
